% !TEX program = lualatex
\documentclass[a4paper,10pt,twocolumn]{article}
\usepackage{kaitoushu}
\renewcommand{\baselinestretch}{1.2}
\lhead{応用数学 問題集}
\problemrange{問74〜80}

\begin{document}
\thispagestyle{fancy}

\begin{center}
  {\Large \textbf{応用数学　2章　第1回　Basic}}
\end{center}
\vspace{0.5em}

\noindent\textbf{\large【Basic】}
\vspace{0.3em}

\mondai{74}{
$f(t) = t^3$ のラプラス変換を求めよ．
\hfill {\scriptsize [教p.47(問1)]}
}

\mondai{75}{
ラプラス変換の線形性を用いて，$\mathcal{L}[3t^2 + 2]$ を求めよ．
\hfill {\scriptsize [教p.48(問2)]}
}

\mondai{76}{
$\mathcal{L}[e^t - e^{-2t}]$ を求めよ．
\hfill {\scriptsize [教p.48(問3)]}
}

\mondai{77}{
定義に従って，$f(t) = \sin 3t$ のラプラス変換を求めよ．
\hfill {\scriptsize [教p.49(問4)]}
}

\mondai{78}{
$f(t) = \cosh 2t$ のラプラス変換を求めよ．
\hfill {\scriptsize [教p.49(問5)]}
}

\mondai{79}{
次の関数のグラフをかけ．また，そのラプラス変換を求めよ．
(1) $y = U(t-3)$ \quad (2) $y = 3U(t-2)$
\hfill {\scriptsize [教p.49(問6)]}
}

\mondai{80}{
次の関数 $f(t)$ を単位ステップ関数を用いて表せ．また，$\mathcal{L}[f(t)]$ を求めよ．
(1) $f(t) = \begin{cases} 1 & (1 \leq t < 3) \\ 0 & (0 < t < 1,\ t \geq 3) \end{cases}$
(2) $f(t) = \begin{cases} 2 & (t \geq 1) \\ 0 & (0 < t < 1) \end{cases}$
\hfill {\scriptsize [教p.50(問7)]}
}

\end{document}